% Announcement note for the LeanVOA library and the book "Vertex Operator Algebras in Lean".
\documentclass[10pt]{article}
\usepackage[T1]{fontenc}
\usepackage[utf8]{inputenc}
\usepackage{lmodern}
\usepackage[margin=0.9in]{geometry}
\usepackage{amsmath,amssymb}
\usepackage{array,booktabs,longtable}
\usepackage{xcolor}
\usepackage{url}
\usepackage[colorlinks=true,linkcolor=blue!50!black,citecolor=blue!50!black,urlcolor=blue!50!black]{hyperref}

% Lean declaration names: typewriter, breakable at dots and underscores.
\DeclareUrlCommand\lean{\urlstyle{tt}}
\makeatletter
\g@addto@macro{\UrlBreaks}{\do\.\do\_}
\makeatother

\setlength{\emergencystretch}{2em}
\newcommand{\todo}[1]{\textcolor{red}{\textbf{[#1]}}}
\newcommand{\C}{\mathbb{C}}
\newcommand{\Z}{\mathbb{Z}}
\newcommand{\sltwo}{\mathfrak{sl}_2}
\newcommand{\slthree}{\mathfrak{sl}_3}
\newcommand{\g}{\mathfrak{g}}

\title{Vertex Operator Algebras in Lean:\\ a formal library and its companion book}
\author{Fedor K.\ Popov\textsuperscript{a} \quad Leonardo Rastelli\textsuperscript{a,b}\\[0.8em]
  \small\textsuperscript{a}Simons Center for Geometry and Physics, Stony Brook University, Stony Brook, NY 11794-3840, USA\\
  \small\textsuperscript{b}C.\,N.\ Yang Institute for Theoretical Physics, Stony Brook University, Stony Brook, NY
  11794-3840, USA\\[0.5em]
  \small\texttt{fpopov@scgp.stonybrook.edu}, \texttt{leonardo.rastelli@stonybrook.edu}}
\date{November 3, 2026}

\begin{document}
\maketitle

\begin{abstract}
  We present \texttt{LeanVOA}, a Lean~4 library for the theory of vertex operator algebras (VOAs), together with a
  companion book that sets each formalized result next to its textbook statement. The library constructs the
  Moonshine module $V^\natural$ and formalizes Zhu's modular invariance, Huang's tensor categories with the Verlinde
  formula, the Gannon--Rayhaun classification of $c=1$ VOAs, and classical freeness of the boundary admissible
  $L_k(\sltwo)$ (Arakawa--Dai) and of the minimal models $M(2,2m+3)$ (van Ekeren--Heluani). It also contains new
  results: graded first-order rigidity of all Virasoro minimal models and of $L_k(\sltwo)$ at every admissible level; classical freeness of $L_{-3/2}(\slthree)$, and of the other members of the Deligne series under stated
  assumptions; and a proof that graded
  unitarity of the boundary admissible $L_k(\sltwo)$ implies that of the minimal models $M(2,2n+1)$. The library
  declares no axiom of its own.
\end{abstract}

% =====================================================================
\section{Introduction}
% =====================================================================

The library \texttt{LeanVOA} formalizes, in Lean~4 \cite{lean4} on top of Mathlib \cite{mathlib}, the algebraic
theory of vertex operator algebras: formal calculus and locality from first principles, the VOA axioms
and the Borcherds identity, the standard constructions (Heisenberg, Virasoro, affine, lattice), the module theory,
Zhu's algebra and modular invariance, Huang's tensor categories, and several recent theorems from the literature
(listed in Section~\ref{sec:scope}). The companion book, \emph{Vertex Operator Algebras in Lean}, is the
reviewed account of the library: for each definition and theorem it gives the mathematical statement, the Lean
declaration copied verbatim from the source, and a ``dictionary'' that decodes every convention and hypothesis
and says whether the Lean statement is the same as, stronger than, or weaker than the textbook one.

\begin{itemize}
  \item Repository: \url{https://github.com/Fedor-Popov/Lean_VOA}
  \item Release: v1.0 (tag \texttt{v1.0})
  \item Book (PDF): \texttt{Notes/book.pdf} in the repository (2{,}337 pages); its \LaTeX{} sources are in
    \texttt{Notes/Textbook/}
  \item Short version of the book (PDF): \texttt{Notes/brief.pdf} (219 pages): the main statements of each chapter,
    each with its verbatim Lean declaration and its status, without proofs or discussion, numbered as in the book
  \item License: Apache-2.0.
\end{itemize}

The library defines VOAs in the general sense of Frenkel--Lepowsky--Meurman and Lepowsky--Li \cite{FLM,LepowskyLi}
(\lean{GradedVertexOperatorAlgebra}), but develops the theory for VOAs of CFT type ($V_0=\C\mathbf 1$, $V_n=0$ for
$n<0$; \lean{VertexOperatorAlgebra}). Every theorem below about VOAs is about VOAs of CFT type; extending the theory
to the general notion is future work.

% =====================================================================
\section{Scope}\label{sec:scope}
% =====================================================================

Table~\ref{tab:scope} lists the headline results: classical theorems, recent theorems from the literature, and
new results. Each entry says what is proved and what, if anything, it assumes; the book gives the precise statements,
conventions and caveats. Every listed declaration depends only on Lean's three standard axioms (\lean{propext},
\lean{Classical.choice}, \lean{Quot.sound}).

{\footnotesize
  \setlength{\tabcolsep}{4pt}
  \renewcommand{\arraystretch}{1.15}
  \begin{longtable}{>{\raggedright\arraybackslash}p{0.625\textwidth}>{\raggedright\arraybackslash}p{0.33\textwidth}}
    \caption{Headline results}\label{tab:scope}\\
    \toprule
    Statement and source & Lean name(s)\\
    \midrule
    \endfirsthead
    \toprule
    Statement and source & Lean name(s)\\
    \midrule
    \endhead
    \bottomrule
    \endfoot
    \multicolumn{2}{l}{\textbf{Tier 1: foundations and classical theorems}}\\
    \midrule
    \textbf{Foundations.} VOAs of CFT type, with the Borcherds identity derived from locality rather than assumed.
    Constructed examples: the Heisenberg VOA, the universal Virasoro VOA at every central charge, the universal affine
    VOA $V^k(\sltwo)$, and the lattice VOA of every positive-definite even lattice. Source: \cite{FLM,LepowskyLi}.
    & \lean{VOA.VertexOperatorAlgebra.borcherds}, \lean{VOA.heisenbergVOA}, \lean{VOA.VirModule.virasoroVOA},
    \lean{VOA.AffineVOA.affineSl2VOA_closed}, \lean{VOA.LatGen.latticeVOAGen}\\
    \addlinespace
    \textbf{Moonshine module.} The Frenkel--Lepowsky--Meurman module $V^\natural$, with no hypotheses: a simple VOA
    with $c=24$, $V^\natural_1=0$ and $\dim V^\natural_2=196884$, whose graded dimensions are the coefficients of
    $j(\tau)-744$. No Monster action is claimed. Source: \cite{FLM,DGM}.
    & \lean{VOA.Moonshine.moonshine}, \lean{VOA.Moonshine.hasSum_JShift}\\
    \addlinespace
    \textbf{Zhu's modular invariance.} For a $C_2$-cofinite rational VOA, the torus one-point functions of the
    irreducible modules transform according to a representation of $SL_2(\Z)$. Instances with every hypothesis checked: all
    rank-one lattice VOAs and their $\Z_2$-orbifolds. Source: \cite{Zhu,DLMtwisted,DLMalg}.
    & \lean{VOA.zhuModularity_canonical_of_C2Cofinite_admissiblyRational},
    \lean{VOA.Classification.COne.DLMRat.admissiblyRational_of_dlmRational},
    \lean{VOA.Classification.COne.zhuModularity_V2m}\\
    \addlinespace
    \textbf{Affine VOAs for every simple Lie algebra.} The universal affine VOA $V^k(\g)$ and its simple quotient
    $L_k(\g)$ at every non-critical level, with concrete models of $\mathfrak{sl}_n$ and of the exceptional Lie
    algebras. Source: \cite{Kac}.
    & \lean{VOA.AffineGeneral.universalAffineVOA}, \lean{VOA.AffineGeneral.simpleAffineQuotientVOA_isSimple}\\
    \midrule
    \multicolumn{2}{l}{\textbf{Tier 2: recent theorems}}\\
    \midrule
    \textbf{Classical freeness of $L_k(\sltwo)$ (Arakawa--Dai).} At the boundary admissible levels $k=-2+2/(2m+1)$, $m\ge1$,
    $L_k(\sltwo)$ is classically free, with $C_2$ algebra $\C[e,h,f]/(eQ^m,hQ^m,fQ^m)$, $Q$ the Casimir; also its vacuum
    character, in product and fermionic form. Source: \cite{ArakawaDai}.
    & \lean{VOA.AffineSl2BoundaryClassicalFreeness.arakawaDai_theorem_1_1},
    \lean{VOA.AffineSl2BoundaryCharacter.character_eq_fermionic_chain}\\
    \addlinespace
    \textbf{Classical freeness of the minimal models $M(2,2m+3)$.} For every $m\ge0$, with $C_2$ algebra
    $\C[T]/(T^{m+1})$. Source: \cite{vanEkerenHeluani}.
    & \lean{VOA.VirasoroClassicalFreeness.classicallyFree}\\
    \addlinespace
    \textbf{Huang's tensor categories.} For strongly rational VOAs, the braided tensor category of modules; for
    $C_2$-cofinite VOAs, that of logarithmic modules. If the VOA is moreover pseudo-unitary: rigidity, the Verlinde
    formula and the congruence property of its modular representation.
    Source: \cite{HuangRigidity,HuangVerlinde,HuangLog,DongLinNg}.
    & \lean{VOA.Huang.TensorCat.HexagonClose.braidedCategory_stronglyRational_pbw},
    \lean{VOA.Huang.HuangRigid.rigid_ordinary},
    \lean{VOA.Classification.COne.huangVerlindeFusionInput_holds},
    \lean{VOA.Huang.CongruenceNamed.dlnCongruenceInput_holds}\\
    \addlinespace
    \textbf{Classification at $c=1$ (Gannon--Rayhaun).} The strongly rational, pseudo-unitary VOAs of central charge one
    are exactly the lattice VOAs
    $V_{2m}$, their orbifolds $V_{2m}^+$, and the three exceptional orbifolds $V_2^{A_4}$, $V_2^{S_4}$, $V_2^{A_5}$, each
    occurring once. Source: \cite{GannonRayhaun}.
    & \lean{VOA.Classification.COne.gannonRayhaun_classification}\\
    \addlinespace
    \textbf{Modular differential equations beyond $C_2$-cofiniteness.} $L_{-4/3}(\sltwo)$ is not $C_2$-cofinite, yet its
    characters converge, its torus conformal blocks form a finite-dimensional space, and its vacuum one-point functions
    satisfy a modular linear differential equation. The Arakawa--Kawasetsu theorem is proved for every VOA with
    finite-dimensional zeroth Poisson homology, the property of quasi-lisse VOAs that it uses. Source:
    \cite{ArakawaKawasetsu}.
    & \lean{VOA.Sl2NonLisse.not_C2Cofinite_simpleSl2QL},
    \lean{VOA.IsSqPrimary.exists_monicMLDEOn_all_analyticBlocks_of_etingofFinite_uncond}\\
    \midrule
    \multicolumn{2}{l}{\textbf{New results: rigidity}}\\
    \midrule
    \textbf{Virasoro minimal models.} Every graded first-order deformation of every minimal model $L(c_{p,q},0)$,
    $p,q\ge2$ coprime, is trivial. For the series $M(2,2n+1)$, rigidity was also proved by Butson and Nair by other
    methods \cite{ButsonNair}. The universal Virasoro VOA is proved not rigid, so the notion has content.
    & \lean{VOA.VirModule.minimal_graded_rigidity}, \lean{VOA.VirPoly.not_isRigidGraded_virasoroVOA}\\
    \addlinespace
    \textbf{Affine $\sltwo$ at admissible levels.} Every graded first-order deformation of $L_k(\sltwo)$ is trivial, at
    every admissible level. This contains Linshaw and Qi's results for $\sltwo$, at the positive integral levels and at
    $k=-4/3$ \cite{LinshawQi}.
    & \lean{VOA.AffineSl2GeneralRigidity.admissible_graded_rigidity}\\
    \midrule
    \multicolumn{2}{l}{\textbf{New results: classical freeness for the Deligne series}}\\
    \midrule
    \textbf{$A_2$, unconditional.} $L_{-3/2}(\slthree)$ is classically free: its $C_2$ algebra is the quotient by eight
    explicit quadrics, and the arc algebra of that quotient is isomorphic to the associated graded of the VOA.
    & \lean{VOA.DeligneAudit.DeligneSeries.A2_classicallyFree_final}\\
    \addlinespace
    \textbf{$D_4$, conditional on one published result.} The same for $L_{-2}(\mathfrak{so}_8)$, with $105$ quadrics,
    assuming a published dimension bound for certain fusion products of current-algebra modules \cite{KusLittelmann}.
    One step of its published proof needs repair; nothing suggests the bound is false.
    & \lean{VOA.DeligneAudit.DeligneSeries.D4_classicallyFree_of_inputT}\\
    \addlinespace
    \textbf{$G_2$, $F_4$, $E_6$, $E_7$, $E_8$, conditional on a conjecture.} The same at the other levels of the series,
    assuming the same published bound, a published result on Demazure modules \cite{Ravinder,ChariVenkatesh,KumarDCF,MathieuDCF}
    (proved in Lean for $D_4$, not yet for these types), and a character identity that we checked numerically to high
    order but that has no published proof.
    & \lean{VOA.DeligneAudit.DeligneSeries.G2_classicallyFree_final} and its $F_4$, $E_6$, $E_7$, $E_8$ analogues\\
    \midrule
    \multicolumn{2}{l}{\textbf{New results: graded unitarity}}\\
    \midrule
    \textbf{Universal affine $\sltwo$, unconditional.} For every real $k<-2$, $V^k(\sltwo)$ is graded unitary in the
    sense of \cite{ABLR}.
    & \lean{VOA.AffineSl2BelowCritical.universal_gradedUnitary_belowCritical_unconditional}\\
    \addlinespace
    \textbf{Transfer through the BRST coset.} If $L_k(\sltwo)$ at the boundary admissible level is graded unitary,
    then so is the minimal model $M(2,2n+1)$, for every $n\ge2$. This is a proved implication; its hypothesis is open.
    The proof formalizes the argument of Beem and Garner \cite{BG} on the actual BRST complex.
    & \lean{VOA.TwoSeries.JetPoincareTwo.twoSeries_gradedUnitary_of_affine_final}\\
  \end{longtable}
}

\paragraph{Also in the library.} Vertex operator superalgebras, with symplectic fermions and the $\beta\gamma$
system as examples, and the triplet algebra $W(2)$, proved simple, $C_2$-cofinite and not rational.

% =====================================================================
\section{Trust model}\label{sec:trust}
% =====================================================================

\paragraph{Base.} The library is built with Lean~4.30.0 and Mathlib v4.30.0. It declares no axiom of its own.
Production code contains no \lean{sorry}, \lean{admit}, \lean{native_decide}, \lean{stop} or
\lean{implemented_by}. We wrote an independent repository checker that rejects any of these, as well as any
\lean{axiom} declaration, and requires every library file to be reachable from the umbrella module. Large finite
verifications are done by kernel reduction (\lean{decide}, \texttt{decide +kernel}), never by the compiler. In addition,
we run \texttt{\#print axioms} on every headline declaration.

\paragraph{Definitions and theorems.} The real risk in a formalization of this kind is not an unsound proof but a true
theorem about the wrong object: a definition that is weaker than the textbook one, a structure that can only be
trivial, a conclusion satisfied by a degenerate witness.
The book is our main defense: it shows how each declaration is connected to the corresponding statement in the
literature.
For every theorem it prints the Lean declaration verbatim (a script checks that each
quoted block still occurs in the source at the recorded commit), decodes every convention and hypothesis in a
dictionary box, and states how the statement could have been wrong or vacuous and why it is not. Each chapter has
two referees: one for clarity, one for fidelity (statement against Lean, verbatim quotes). Before release, each
headline area received an independent \emph{release audit}: a fresh AI reviewer and human readers
compared the Lean statement with the source paper, traced every definition it depends on,
listed every explicit hypothesis, and re-printed the axioms. A final read of the book's release edition then brought every chapter into line with the audited wording.

% =====================================================================
\section{Size and structure}
% =====================================================================

The released library has 11{,}660 Lean files and 3{,}752{,}244 lines under \texttt{VOA/}.

\begin{center}\small
  \begin{tabular}{lrrl}
    \toprule
    Directory & Files & Lines & Content\\
    \midrule
    \texttt{Foundations} & 29 & 9{,}642 & formal calculus, PBW, Lie-theoretic input\\
    \texttt{Core} & 128 & 46{,}320 & VOA axioms, Borcherds identity, quotients\\
    \texttt{Modules} & 466 & 132{,}816 & weak/admissible/ordinary modules, Zhu algebra\\
    \texttt{Modularity} & 234 & 66{,}958 & modular forms, Zhu's theorem, $q$-traces\\
    \texttt{Examples} & 1{,}181 & 307{,}367 & constructions; rigidity; Arakawa--Dai; graded unitarity\\
    \texttt{AffineGeneral} & 354 & 21{,}641 & affine VOAs for every simple $\g$\\
    \texttt{LatticesAndCodes} & 25 & 8{,}472 & lattices, Golay code, Leech lattice\\
    \texttt{Moonshine} & 170 & 55{,}246 & $V^\natural$\\
    \texttt{Huang} & 931 & 183{,}719 & tensor categories, rigidity, Verlinde, congruence\\
    \texttt{Classification} & 1{,}756 & 1{,}627{,}513 & Gannon--Rayhaun $c=1$ classification\\
    \texttt{DeligneAudit} & 6{,}173 & 1{,}241{,}760 & Deligne series\\
    \texttt{Commutative}, \texttt{Filtered}, \texttt{Geometric} & 201 & 45{,}436 & commutative algebra, Li filtration\\
    umbrella files & 12 & 5{,}354 & \\
    \bottomrule
  \end{tabular}
\end{center}

A large part of the line count is certificate data. The 167 files larger than 200\,KB contain 1{,}518{,}787 lines
(about 40\% of the total), mostly exact rational data for the $c=1$ orbifold computations; they are checked by kernel
reduction (about 39{,}600 \texttt{decide +kernel} calls), and the book describes them collectively --- what is checked,
how, and at what size --- while presenting the theorems that consume them individually.

The release edition of the book has 2{,}337 pages, and its brief edition 219 pages; every Lean block quoted in the brief
(1{,}315) is checked verbatim against the source.

% =====================================================================
\section{Method}
% =====================================================================

The formalization was human-directed and carried out with the assistance of AI coding agents (Anthropic's
Claude Opus and Claude Fable, and OpenAI's Codex), while Python scripts generated the certificate data that the Lean
kernel re-checks and Python and Mathematica supplied independent numerical cross-checks. The authors chose the
targets, the definitions and conventions, and the release scope; agents wrote all of the Lean code and the first
drafts of the book chapters. Three practices shaped the result. First, referee loops: every chapter and every
headline went through independent review by several open-weight and proprietary AI models, and statements were
corrected or narrowed when a review found an over-claim. Second, a strict reporting rule: a result counts only when the aggregate build is green, and
``unconditional'' is reserved for statements whose hypothesis closure is fully discharged. Third, agents were required to record every dead
end, which let later agents narrow their search for proofs:
routes proved false, vacuous or circular are kept, with their counterexamples, in a project log, and some are shown as
examples in the book's verification chapter.

% =====================================================================
\section{In preparation}\label{sec:prep}
% =====================================================================

\begin{itemize}
  \item \textbf{Deligne series beyond $A_2$.} Removing the assumptions of the conditional theorems; for $D_4$ the
    book (Chapter~L27) describes the current state.
  \item \textbf{Graded unitarity in the other direction}, from the minimal models to the affine algebra, through
    Adamovi\'c's free-field realization (inverse Drinfeld--Sokolov reduction).
  \item \textbf{Affine $\slthree$}, towards its rigidity.
  \item \textbf{Removing pseudo-unitarity.} Proving rigidity, the Verlinde formula and the congruence property for
    every strongly rational VOA, as in \cite{HuangRigidity,HuangVerlinde,DongLinNg}, without assuming
    pseudo-unitarity.
\end{itemize}

% =====================================================================
\section*{Acknowledgements}
% =====================================================================
We thank Brandon Rayhaun, Avia Raviv-Moshe and Chris Beem for useful discussions. We also thank the Lean and
Mathlib communities. The work of LR is supported in part by NSF grant PHY-2513893 and by the Simons Foundation
grant 681267 (Simons Investigator Award).

% =====================================================================
\begin{thebibliography}{99}\footnotesize
  % =====================================================================

  \bibitem{ABLR}
  A.~Arabi Ardehali, C.~Beem, M.~Lemos and L.~Rastelli, Graded unitarity in the SCFT/VOA correspondence,
  arXiv:2507.23781.
  \bibitem{BG}
  C.~Beem and A.~Garner, Graded unitarity and relative semi-infinite cohomology, arXiv:2509.10364.
  \bibitem{FLM}
  I.~Frenkel, J.~Lepowsky and A.~Meurman, \emph{Vertex Operator Algebras and the Monster}, Pure and Applied
  Mathematics~134, Academic Press, 1988.

  \bibitem{LepowskyLi}
  J.~Lepowsky and H.~Li, \emph{Introduction to Vertex Operator Algebras and Their Representations}, Progress in
  Mathematics~227, Birkh\"auser, 2004.

  \bibitem{Kac}
  V.~G.~Kac, \emph{Infinite-Dimensional Lie Algebras}, 3rd ed., Cambridge University Press, 1990.

  \bibitem{Zhu}
  Y.~Zhu, Modular invariance of characters of vertex operator algebras, \emph{J.\ Amer.\ Math.\ Soc.}\ \textbf{9}
  (1996), 237--302.

  \bibitem{DLMtwisted}
  C.~Dong, H.~Li and G.~Mason, Twisted representations of vertex operator algebras, \emph{Math.\ Ann.}\ \textbf{310}
  (1998), 571--600, arXiv:q-alg/9509005 (Theorem~8.1).

  \bibitem{DLMalg}
  C.~Dong, H.~Li and G.~Mason, Vertex operator algebras and associative algebras, \emph{J.\ Algebra} \textbf{206}
  (1998), 67--96 (Theorem~4.10).

  \bibitem{DGM}
  L.~Dolan, P.~Goddard and P.~Montague, Conformal field theories, representations and lattice constructions,
  \emph{Comm.\ Math.\ Phys.}\ \textbf{179} (1996), 61--120, arXiv:hep-th/9410029.

  \bibitem{HuangRigidity}
  Y.-Z.~Huang, Rigidity and modularity of vertex tensor categories, \emph{Commun.\ Contemp.\ Math.}\ \textbf{10}
  (2008), suppl.~1, 871--911, arXiv:math/0502533.

  \bibitem{HuangVerlinde}
  Y.-Z.~Huang, Vertex operator algebras and the Verlinde conjecture, \emph{Commun.\ Contemp.\ Math.}\ \textbf{10}
  (2008), no.~1, 103--154, arXiv:math/0406291.

  \bibitem{HuangLog}
  Y.-Z.~Huang, Cofiniteness conditions, projective covers and the logarithmic tensor product theory, \emph{J.\ Pure
  Appl.\ Algebra} \textbf{213} (2009), 458--475, arXiv:0712.4109; and Y.-Z.~Huang, J.~Lepowsky and L.~Zhang,
  Logarithmic tensor category theory for generalized modules for a conformal vertex algebra, I: Introduction and
  strongly graded algebras and their generalized modules, in \emph{Conformal Field Theories and Tensor Categories},
  Math.\ Lect.\ Peking Univ., Springer, 2014, 169--248, arXiv:1012.4193 (Part~I of a series of eight).

  \bibitem{DongLinNg}
  C.~Dong, X.~Lin and S.-H.~Ng, Congruence property in conformal field theory, \emph{Algebra Number Theory}
  \textbf{9} (2015), no.~9, 2121--2166.

  \bibitem{GannonRayhaun}
  T.~Gannon and B.~C.~Rayhaun, Classification of rational $c=1$ vertex operator algebras and vertex operator
  superalgebras, arXiv:2609.00122 (2026).

  \bibitem{ArakawaDai}
  T.~Arakawa and X.~Dai, Classical freeness of affine vertex algebras of $\sltwo$ at boundary admissible levels,
  arXiv:2609.33348 (2026).

  \bibitem{ArakawaKawasetsu}
  T.~Arakawa and K.~Kawasetsu, Quasi-lisse vertex algebras and modular linear differential equations, in \emph{Lie
  Groups, Geometry, and Representation Theory}, Progr.\ Math.~326, Birkh\"auser, 2018, 41--57, arXiv:1610.05865.

  \bibitem{LinshawQi}
  A.~R.~Linshaw and F.~Qi, Deformation rigidity of some simple affine VOAs, arXiv:2601.12017 (2026).

  \bibitem{ButsonNair}
  D.~Butson and S.~Nair, On the deformation theory of chiral quantizations, arXiv:2606.27341.
  \bibitem{vanEkerenHeluani}
  J.~van~Ekeren and R.~Heluani, Chiral homology of elliptic curves and the Zhu algebra, \emph{Comm.\ Math.\ Phys.}\
  \textbf{386} (2021), 495--550, arXiv:1804.00017.

  \bibitem{Ravinder}
  B.~Ravinder, Demazure modules, Chari--Venkatesh modules and fusion products, \emph{SIGMA} \textbf{10} (2014), 110,
  arXiv:1409.0274.

  \bibitem{KusLittelmann}
  D.~Kus and P.~Littelmann, Fusion products and toroidal algebras, \emph{Pacific J.\ Math.}\ \textbf{278} (2015),
  427--445, arXiv:1411.5272.

  \bibitem{ChariVenkatesh}
  V.~Chari and R.~Venkatesh, Demazure modules, fusion products and $Q$-systems, \emph{Comm.\ Math.\ Phys.}\ \textbf{333}
  (2015), 799--830, arXiv:1305.2523.

  \bibitem{KumarDCF}
  S.~Kumar, Demazure character formula in arbitrary Kac--Moody setting, \emph{Invent.\ Math.}\ \textbf{89} (1987),
  395--423.

  \bibitem{MathieuDCF}
  O.~Mathieu, Formules de caract\`eres pour les alg\`ebres de Kac--Moody g\'en\'erales, \emph{Ast\'erisque}
  \textbf{159--160} (1988).

  \bibitem{lean4}
  L.~de~Moura and S.~Ullrich, The Lean~4 theorem prover and programming language, in \emph{Automated Deduction ---
  CADE~28}, Lecture Notes in Comput.\ Sci.~12699, Springer, 2021, 625--635.

  \bibitem{mathlib}
  The mathlib Community, The Lean mathematical library, in \emph{Proceedings of the 9th ACM SIGPLAN International
  Conference on Certified Programs and Proofs (CPP 2020)}, ACM, 2020, 367--381.

\end{thebibliography}

\end{document}
